\documentclass[aspectratio=169,10pt]{beamer}
\newcommand{\reporoot}{../..}
% Shared preamble of the main deck and the topic decks.
% Define \reporoot (the path from the .tex file to the repository root)
% before \input-ing this file.
\usetheme[progressbar=frametitle,numbering=fraction,sectionpage=progressbar]{metropolis}

\usepackage{amsmath,amssymb}
\usepackage{booktabs}
\usepackage{graphicx}
\usepackage{tikz}
\usetikzlibrary{arrows.meta,positioning,decorations.pathreplacing}
\usepackage[authoryear,round]{natbib}
\usepackage{xcolor}

\graphicspath{{\reporoot/figures/}}

% Every number on the slides is produced by experiments/*.py and read from here.
\input{\reporoot/results/e1_markov_book.tex}
\input{\reporoot/results/e2_best_queues.tex}
\input{\reporoot/results/e3_queue_reactive.tex}
\input{\reporoot/results/e4_hawkes.tex}
\input{\reporoot/results/e5_signature.tex}
\input{\reporoot/results/e6_propagator.tex}
\input{\reporoot/results/e7_market_making.tex}
\input{\reporoot/results/e8_ofi.tex}
\input{\reporoot/results/e9_queue_value.tex}
\input{\reporoot/results/e10_markouts.tex}
\input{\reporoot/results/e11_glft_exact.tex}
\input{\reporoot/results/e12_signal.tex}
\input{\reporoot/results/e13_lvr.tex}
\input{\reporoot/results/e14_mlofi.tex}
\input{\reporoot/results/e15_rate_estimation.tex}
\input{\reporoot/results/e16_queue_reactive_price.tex}
\input{\reporoot/results/e17_hawkes_pitfalls.tex}
\input{\reporoot/results/e18_propagator_estimation.tex}
\input{\reporoot/results/e19_latent_book.tex}
\input{\reporoot/results/e20_microprice.tex}
\input{\reporoot/results/e21_fill_estimation.tex}
\input{\reporoot/results/e22_glosten_milgrom.tex}
\input{\reporoot/results/e23_cdl_on_book.tex}

\newcommand{\E}{\mathbb{E}}
\newcommand{\Var}{\operatorname{Var}}
\newcommand{\OFI}{\mathrm{OFI}}
\newcommand{\dd}{\mathrm{d}}
\newcommand{\src}[1]{\par\vfill{\scriptsize\color{gray}#1}}
\definecolor{accent}{HTML}{0C5DA5}
\definecolor{warn}{HTML}{C0392B}
\newcommand{\hl}[1]{\textcolor{accent}{\textbf{#1}}}

\setbeamerfont{footnote}{size=\tiny}
\setbeamerfont{caption}{size=\scriptsize}

\author{Yaohui Shu}
\institute{Stony Brook University}
\date{October 2026}

% Topic decks: a references frame that lists only the works cited in the deck.
\newcommand{\deckreferences}{%
  \begin{frame}[allowframebreaks]{References}
    \tiny
    \renewcommand{\bibsection}{}
    \bibliographystyle{plainnat}
    \bibliography{\reporoot/slides/refs}
  \end{frame}}

% Topic decks: the standard "where this fits" footer on the title frame.
\newcommand{\deckseries}{Deep dive in the series \emph{Stochastic Models of the Limit Order Book}}

% Topic decks: \decknumber{n} puts the series line under the date on the title page.
\newcommand{\decknumber}[1]{\institute{Stony Brook University\hfill{\scriptsize\color{gray}Deep dive #1 of 9}}}


\title{Self-Exciting Order Flow}
\subtitle{Hawkes processes: model, estimation, and how the branching ratio is misread}
\decknumber{4}

\begin{document}

\maketitle

\begin{frame}{The question}
  \small
  \begin{columns}[T]
    \begin{column}{0.55\textwidth}
      Orders cluster in time: one trade triggers others, cancellations come in bursts. \hl{How much of market activity is triggered by market activity itself?}
      \begin{itemize}
        \item The share of endogenous events, the branching ratio, is a measure of reflexivity.
        \item Near-critical excitation links millisecond order flow to rough volatility over hours.
        \item For an LP, clustering means fills and adverse moves arrive in bursts, not evenly.
      \end{itemize}
    \end{column}
    \begin{column}{0.42\textwidth}
      \begin{alertblock}{In this deck}
        \small
        \begin{enumerate}
          \item The process and its cluster picture.
          \item Exact likelihood and diagnostics.
          \item The Hawkes price model and the signature plot.
          \item Three ways an exponential fit misreads the branching ratio.
          \item The criticality debate.
        \end{enumerate}
      \end{alertblock}
    \end{column}
  \end{columns}
\end{frame}

\begin{frame}{The process}
  \small
  For event types $u, v$ (buy, sell, cancel, \ldots) with event times $t^v_j$, the intensity (s$^{-1}$) is
  \[ \Lambda_u(t) = \Lambda^0_u + \sum_v \sum_{t^v_j < t} \phi_{uv}(t - t^v_j), \]
  with baseline $\Lambda^0_u$ (s$^{-1}$) and kernels $\phi_{uv} \geq 0$ (s$^{-1}$). Write $G_{uv} = \int_0^\infty \phi_{uv}(t)\,\dd t$ (dimensionless).
  \begin{columns}[T]
    \begin{column}{0.5\textwidth}
      \textbf{Cluster picture.} Immigrants arrive at rate $\Lambda^0$; each type-$v$ event has on average $G_{uv}$ direct children of type $u$. The process is stationary iff the spectral radius $\rho(G) < 1$, and then
      \[ \bar\Lambda = (I - G)^{-1}\Lambda^0 . \]
      In one dimension $n = G$ is the \hl{branching ratio}: the share of events that are endogenous.
    \end{column}
    \begin{column}{0.46\textwidth}
      \textbf{Exponential kernel} $\phi(t) = \alpha e^{-\beta t}$, with jump $\alpha$ and decay $\beta$ (s$^{-1}$): $n = \alpha/\beta$, and the intensity is Markov, so the likelihood takes one pass.

      \vspace{0.3em}
      \textbf{Power-law kernel} $\phi(t) \propto (t + \tau_0)^{-(1+\epsilon)}$, with cut-off $\tau_0$ (s) and tail exponent $\epsilon > 0$: memory over many time scales, as in data.
    \end{column}
  \end{columns}
\end{frame}

\begin{frame}{Likelihood and diagnostics}
  \small
  On $[0, T]$ the log-likelihood of a point process with intensity $\Lambda_u$ is
  \[ \ell = \sum_u\Big[\sum_j \ln\Lambda_u(t^u_j) - \int_0^T \Lambda_u(t)\,\dd t\Big]. \]
  For exponential kernels both terms update recursively, so $\ell$ costs $O(N)$ for $N$ events. Standard errors come from the inverse observed information.
  \vspace{0.3em}
  \begin{columns}[T]
    \begin{column}{0.5\textwidth}
      \textbf{Time-rescaling test.} With compensator $\mathcal{A}_u(t) = \int_0^t \Lambda_u$, the increments $\mathcal{A}_u(t^u_j) - \mathcal{A}_u(t^u_{j-1})$ are i.i.d.\ unit exponential under the true model. A Kolmogorov-Smirnov (KS) test checks them.
    \end{column}
    \begin{column}{0.46\textwidth}
      \textbf{Nonparametric kernels.} For slowly decaying kernels, \citet{bacry2016estimation} estimate $\phi$ from the empirical covariance of event counts by solving a Wiener-Hopf equation, over lags from microseconds to minutes.
    \end{column}
  \end{columns}
\end{frame}

\begin{frame}{Fitting buy and sell market orders}
  \centering
  \includegraphics[width=0.8\textwidth]{fig4_hawkes.png}
  \begin{columns}[T]
    \begin{column}{0.48\textwidth}
      \footnotesize
      Bivariate exponential Hawkes model of buys and sells, $T = \EfourT$ s: $\EfourNbuy$ buys and $\EfourNsell$ sells, true spectral radius $\rho(G) = \EfourRho$. (a) Intensities over one minute.
    \end{column}
    \begin{column}{0.48\textwidth}
      \footnotesize
      MLE recovers all ten parameters within $\EfourMaxZ$ standard errors. (b) Rescaled waiting times are unit exponential: KS $p = \EfourKSbuy$ and $\EfourKSsell$.
    \end{column}
  \end{columns}
\end{frame}

\begin{frame}{A Hawkes price: microstructure noise from cross-excitation}
  \begin{columns}[c]
    \begin{column}{0.5\textwidth}
      \small
      Up ticks $N_+$ and down ticks $N_-$, price $X = N_+ - N_-$ (ticks), self jump $\alpha_s$, cross jump $\alpha_c$, common decay $\beta$ \citep{bacry2013modelling}. With $g = (\alpha_s - \alpha_c)/\beta$ and $\eta_H = \beta(1 - g)$ (s$^{-1}$),
      \[ C(h) = 2\bar\Lambda\Big[\tfrac{1}{(1-g)^2} + \Big(1 - \tfrac{1}{(1-g)^2}\Big)\tfrac{1 - e^{-\eta_H h}}{\eta_H h}\Big] \]
      is the realized variance per unit time at sampling interval $h$ (ticks$^2$ s$^{-1}$), with $\bar\Lambda$ the mean rate per side.
      \begin{itemize}
        \item $\alpha_c > \alpha_s$: an up tick triggers down ticks, $C(h)$ falls with $h$: the signature plot of bid-ask bounce.
        \item Simulation within $\EfiveMaxZ$ s.e.\ over 42 points.
      \end{itemize}
    \end{column}
    \begin{column}{0.46\textwidth}
      \centering
      \includegraphics[width=\textwidth]{fig5_signature.png}
    \end{column}
  \end{columns}
\end{frame}

\begin{frame}{Three ways to misread the branching ratio}
  \centering
  \includegraphics[width=0.92\textwidth]{fig17_hawkes_pitfalls.png}
  \footnotesize
  \begin{columns}[T]
    \begin{column}{0.32\textwidth}
      (a) Truth: a sum of 16 exponentials that follows $(t + \tau_0)^{-1.15}$ from $0.1$ s to about $10^4$ s, with $n = 0.9$.
    \end{column}
    \begin{column}{0.32\textwidth}
      (b) An exponential kernel fitted on windows from 1 min to 8 h finds $\hat n = \EseventeenShort$ to $\EseventeenLong$, with decay time $1/\hat\beta \approx \EseventeenBetaInv$ s. It misses more than half the excitation at every window length.
    \end{column}
    \begin{column}{0.32\textwidth}
      (c) Pure Poisson data ($n = 0$), $\EseventeenDays$ days each: a U-shaped intraday rate gives $\hat n = \EseventeenSeasonRaw$; after a time change by the estimated profile, $\EseventeenSeasonDes$. Random rate regimes: $\EseventeenSwitch$.
    \end{column}
  \end{columns}
\end{frame}

\begin{frame}{Why the fits go wrong}
  \small
  \begin{columns}[T]
    \begin{column}{0.5\textwidth}
      \textbf{Nonstationary baseline.} A slow change in the baseline rate looks to the likelihood like a long cluster. A slowly decaying exponential ($1/\hat\beta$ of about 100 s here) absorbs it and reports $\hat n$ close to one. This is the result of \citet{filimonov2015apparent}: Poisson data with regime changes give spurious near-critical branching ratios.

      \vspace{0.3em}
      \textbf{Remedy.} Estimate the seasonal profile across days and change time by its integral, $t \mapsto \int_0^t \hat\Lambda^0(s)\,\dd s$, before fitting. Regime changes need a model of their own.
    \end{column}
    \begin{column}{0.46\textwidth}
      \textbf{Misspecified kernel.} A single exponential has one time scale. Faced with excitation spread over six decades, the likelihood settles on the fast part and reports less than half of $n$, whatever the window.

      \vspace{0.3em}
      \textbf{Remedy.} Use kernels with many time scales (sums of exponentials or power laws), nonparametric estimation, and check residuals with the time-rescaling test.

      \vspace{0.3em}
      \textbf{Other traps.} Coarse timestamps and simultaneous events bias short-lag excitation; edge effects (no history before the window) bias $\hat n$ down.
    \end{column}
  \end{columns}
\end{frame}

\begin{frame}{What the data show}
  \small
  \begin{itemize}
    \item \textbf{Critical reflexivity.} For mid-price changes of the E-mini S\&P 500 futures, 1998 to 2011, a power-law kernel with exponent near $-1.15$ below $10^3$ s and near $-1.45$ from $10^3$ to $10^6$ s integrates to about one in every two-month period: the market is, and was, close to critical \citep{hardiman2013critical}.
    \item \textbf{Why exponential fits disagreed.} Fitting exponential kernels on 30-minute windows to synthetic critical power-law data gives sub-critical branching ratios that rise as the short-time cut-off falls, which can mimic a rise in reflexivity over the years \citep{hardiman2013critical}.
    \item \textbf{The counterpoint.} Regime shifts and outliers bias $\hat n$ upward, and calibration on Poisson data with regime changes gives $n \approx 1$ when the truth is zero \citep{filimonov2015apparent}.
    \item \textbf{Many dimensions.} Nonparametric estimation on an eight-dimensional order book model finds slowly decaying kernels over lags from about 10 $\mu$s to 100 s \citep{bacry2016estimation}. Survey: \citet{bacry2015hawkes}.
  \end{itemize}
\end{frame}

\begin{frame}{From order flow to rough volatility}
  \small
  \begin{itemize}
    \item As $\rho(G) \to 1$ with the horizon, suitably rescaled nearly unstable Hawkes processes converge to an integrated Cox-Ingersoll-Ross-type process \citep{jaisson2015limit}. Later work with heavy-tailed kernels obtains \hl{rough} limits, with Hurst exponent below one half.
    \item Adding no-arbitrage, buy-sell asymmetry and metaorder splitting gives the rough Heston model with leverage \citep{eleuch2018microstructural}; recent work adds common price and volatility jumps \citep{wang2026rough}.
    \item For an LP: the branching ratio estimated from milliseconds of order flow and the volatility hedged over hours are two views of one process. Mis-estimating $n$ (previous slides) mis-estimates both.
  \end{itemize}
  \vspace{0.4em}
  \begin{alertblock}{Takeaways}
    \small
    Fit with kernels that have many time scales; remove seasonality first; test residuals; and distrust a branching ratio near one until regime changes are ruled out. The criticality question is open because each answer depends on these choices.
  \end{alertblock}
\end{frame}

\begin{frame}[fragile]{Reproduce}
  \small
  Code: \texttt{hawkes.py}. Experiments: \texttt{e4\_hawkes.py}, \texttt{e5\_signature.py}, \texttt{e17\_hawkes\_pitfalls.py}.
\begin{verbatim}
from stochastic_lob import hawkes as hk
w, beta = hk.power_law_weights(n=0.9, eps=0.15, tau0=0.1, beta_min=1e-4)
model = hk.SumExpHawkes1D(mu=0.05, w=w, beta=beta)
t = model.simulate(T=3e4, rng=1)
mu, n_hat, beta_hat, se = hk.fit_exp_hawkes_1d(t, 3e4)
\end{verbatim}
\end{frame}

\deckreferences
\end{document}
